% Included after datatype-reference-appendix.tex in the paper.
\section{Five Bugs in Notion's Public Demo}
\label{app:notion-bugs}

The public educational demo accompanying Notion's CRDT blog
post~\cite{NotionCRDTBlog,NotionCRDTDemo} exhibits five failures of ordinary
editing behaviour. Each trace below uses requests generated by the
demo's own client, starting from an empty document. Concurrent requests are
prepared at different replicas with distinct session identities, and a replica
receives an update only after the updates its issuer had seen. In the first
three traces, after receiving the concurrent updates, a replica issues the
final request against its displayed document. The unexpected result is thus
the immediate effect of a local edit, rather than a disagreement between
replicas. In the last two traces, text disappears when a delayed update
arrives, although no request deletes text. The replicas still converge, but
no sequential order of the requests removes a character, so no sequential
execution explains the converged display.

We write $[w_1,\ldots,w_n]$ for the displayed blocks, $\epsilon$ for an
empty block, and bold characters for bold formatting. Positions are
zero-based. Splitting a block $w$ at position $k$ should replace it by
$[\operatorname{take}_k(w),\operatorname{drop}_k(w)]$.
Merging a block into its predecessor appends its text to that predecessor
and removes the block; this editing request is distinct from merging
replicated versions. All characters in the examples are initially plain.
Toggling bold on such a selection should bold precisely the selected
characters and preserve the remaining text and formatting.

\subsection{Splitting at Zero Moves Only Part of the Text}
\label{app:notion-split-bug}

Start with one block containing $a$. Split it at zero and insert $X$ into
the empty first block, obtaining $[X,a]$. Two replicas independently merge
the second block into its predecessor. Deliver both updates to one replica;
its display is $[Xa]$. From this state, two further replicas split the block
concurrently, one at zero and the other at one. After receiving both
updates, merge the block created by the split at zero into its predecessor.
The resulting display is $[\epsilon,aX]$.

Now split the block $aX$ at zero. The expected and observed results are
\[
\begin{array}{rcl}
\text{expected}&:&[\epsilon,\epsilon,aX],\\
\text{observed}&:&[\epsilon,X,a].
\end{array}
\]
The split moves a subtree containing $a$, but leaves a sibling branch
containing $X$ in the original block. Thus it moves only part of the
displayed suffix and reverses the displayed character order. The two earlier
merge-block requests have distinct identifiers: this failure is not caused
by delivering the same request twice.

\subsection{Bolding a Selection Changes Neither Selected Character}
\label{app:notion-empty-bold-bug}

Start with one block containing $ab$. Two replicas split it concurrently
at positions zero and one. Receiving both updates produces
$[\epsilon,b,a]$. Merge the block containing $a$ into its predecessor,
obtaining $[\epsilon,ba]$. Toggle bold on both displayed characters of $ba$:
\[
\begin{array}{rcl}
\text{expected}&:&[\epsilon,\mathbf{ba}],\\
\text{observed}&:&[\epsilon,ba].
\end{array}
\]
Neither character becomes bold. Although the displayed order is now $ba$,
both characters belong to an underlying text instance ordered $ab$.
The demo represents the selection as an interval from immediately before
its first character to immediately after its last character in that
underlying order. Here these boundaries are before $b$ and after $a$:
the same position in $ab$. The interval is empty.

\subsection{Bolding a Block Also Changes Another Block}
\label{app:notion-excess-bold-bug}

Start with one block containing $abc$. Two replicas split it concurrently
at positions one and two. Receiving both updates produces $[a,c,b]$.
Merge the block containing $c$ into its predecessor, obtaining $[ac,b]$.
Toggle bold on the two characters in the first block:
\[
\begin{array}{rcl}
\text{expected}&:&[\mathbf{ac},b],\\
\text{observed}&:&[\mathbf{ac},\mathbf b].
\end{array}
\]
The unselected $b$ also becomes bold. The selected characters $a$ and $c$
are consecutive in the displayed block, but the underlying text instance
still has order $abc$. Its interval from before $a$ to after $c$ includes
$b$, even though $b$ is displayed in a different block.

\subsection{Merging Blocks Hides the Whole Document}
\label{app:notion-document-loss-bug}

Start with three blocks $[a,b,c]$ at two replicas. The first replica merges
$c$ into its predecessor ($[a,bc]$), splits that block at zero
($[a,\epsilon,bc]$), and merges $bc$ back into the empty block ($[a,bc]$).
Concurrently, the second replica merges $b$ into its predecessor ($[ab,c]$).
It then receives the first replica's first two updates, giving
$[\epsilon,abc]$, and merges $abc$ into its predecessor, giving $[abc]$.
Finally, it receives the first replica's last merge:
\[
\begin{array}{rcl}
\text{expected}&:&[abc],\\
\text{observed}&:&[\,].
\end{array}
\]
No block remains visible. The characters are still stored, but only in
removed blocks, and both replicas converge on the empty display. A merge
moves the text after the last text slice of the target block, and locates
that slice by an identifier for its end rather than by block. When the last
merge arrives at the second replica, that slice lies in a block that a merge
has already removed. The demo does not check whether the block is still
present, so the text moves into the removed block.

\subsection{A Delayed Merge Hides a Character}
\label{app:notion-character-loss-bug}

Start with one block containing $x$ at two replicas. The first replica splits
it at zero, giving $[\epsilon,x]$. Concurrently, the second replica splits it
at one ($[x,\epsilon]$) and merges the new empty block back ($[x]$). The
first replica receives only the second replica's split, giving
$[\epsilon,\epsilon,x]$, and merges $x$ into its predecessor, giving
$[\epsilon,x]$. The second replica's merge then arrives:
\[
\begin{array}{rcl}
\text{expected}&:&[\epsilon,x],\\
\text{observed}&:&[\epsilon].
\end{array}
\]
The character $x$ disappears. In the first replica's state, the delayed
merge would move a text slice to a position inside itself. The demo treats
such a move as a no-op, but still removes the merged block, which now
contains $x$. If the merge arrives before the first replica's own merge,
or that merge is not issued, $x$ remains visible.

\paragraph{Reproduction and formal checks.}
We record the source-file hashes and complete request traces in the code repository.
The regression tests reproduce all five outcomes and include nearby cases
that agree with the specification. In Lean, we check that
the split outcome violates the ordinary split-at-zero rule, and that the
interval calculation in the two formatting examples disagrees with
formatting exactly the selected characters. For the last two bugs, we model
a merge's move and block removal over the demo's slice trees, check that
the update hides the text while it stays stored, and prove that every
sequential history of splits and merges preserves the document's characters.
These checks do not formalise
the complete JavaScript execution or prove its correspondence to Lean.
The results concern the public demo; we have not established that Notion's
production system has these bugs. The setups use merge-block and formatting
requests outside the verified create/insert/erase/split fragment, so they
do not contradict its correctness result.
